%%% ВЫБОР УСТРОЙСТВА %%%

\IfStrEqCase{\devicecode}{%
{ЮТКБ.656122.611-06.01 РЭ}{%ДЗТ
\def\fileExt{t}
\def\AbzazOne{
\item
Для исключения ложного срабатывания устройства дуговой защиты, которое может возникнуть при неисправности ее датчиков, предусмотрен контроль по току. Обеспечивается контроль тока трех фаз сторон ВН, НН1(НН), НН2(СН) и токов нейтрали трансформатора. Пороги по току срабатывания индивидуальны для каждой из контролируемых цепей.
}
}%
{ЮТКБ.656122.611-06.02 РЭ}{%ДЗАТ
\def\fileExt{at}
\def\AbzazOne{
\item
Для исключения ложного срабатывания устройства дуговой защиты, которое может возникнуть при неисправности ее датчиков, предусмотрен контроль по току. Обеспечивается контроль тока трех фаз сторон ВН, СН, НН. Пороги по току срабатывания индивидуальны для каждой из контролируемых цепей.
}
}%
}
[
\def\fileExt{??}
\def\AbzazOne{??}
]

\color{unimain}{\subsection{Токовый контроль защиты от дуговых замыканий}}
\color{black}

\begin{enumerate}

\AbzazOne

\item
Алгоритм работы ТК ЗДЗ приведен на рисунке \ref{pic:ZDZ}. Параметры для настройки ТК ЗДЗ приведены в таблице \ref{app_set:tkzdz}.

\medskip
\begin{figure}[H]
\centering
\includegraphics[width=0.75\textwidth,height=0.75\textheight,keepaspectratio]{\fbpath/06.02 ДЗАТ/041.1191 ТК ЗДЗ/img/ZDZ_\fileExt.pdf}
\captionof{figure}{Функциональная схема логики ТК ЗДЗ}\label{pic:ZDZ}
\end{figure}


\end{enumerate}